% =========================================================================
% Section 1 — Introduction (Korean translation)
% File: sections/1_introduction.tex  (\input fragment; no \documentclass)
% Model name: DefectSegAM (exact spelling, everywhere)
% =========================================================================

\section{서론}
\label{sec:intro}

금속 적층제조(additive manufacturing, AM)는 시제품 제작 단계를 넘어 본격적인
양산 단계로 확고히 이동하고 있다. 글로벌 AM 시장은 2023년 200.35억~USD 규모에
도달하여 전년 대비 11.1\% 성장하였고, 금속 AM 장비는 2022년 3{,}049대에서
2023년 3{,}793대로 판매량이 증가하였다~\cite{Wohlers2024}. 이를 보완하는 산업
분석에서는 산업용 AM 시장을 2023년 105.0억~EUR 규모로 평가하고 2028년까지
연평균 13.9\%의 성장률을 전망하였으며, 금속 AM이 폴리머 AM보다 약 두 배
빠르게 성장할 것으로 예측하였다~\cite{AMPOWER2024}. 레이저 분말 베드
융합(laser powder bed fusion, L-PBF; ISO/ASTM 52900 표기로는
PBF-LB/M)~\cite{ISO52900,Gibson2021AMTech,King2015LPBF}은
항공우주 및 생체의료 분야에서 높은 가치를 갖는 기하학적 자유도와 함께
근사 최종 형상(near-net-shape)의 금속 부품을 제작할 수 있기 때문에 이러한
성장의 상당 부분을 차지한다~\cite{BlakeyMilner2021Aerospace,Bose2018Biomaterials}.
그러나 동일한 층별 용융 공정은 키홀 기공(keyhole
porosity)~\cite{King2014Keyhole}, 용융 부족(lack of
fusion)~\cite{Tang2017LOF}, 스패터(spatter) 비산~\cite{Ly2017Spatter}, 그리고
리코터--파트 간섭과 같은 확률적 결함 형성에 취약하며, 이러한 결함의 발생은
국소 형상과 공정 이력에 민감하게 결합되어
있다~\cite{Sames2016Metallurgy,DebRoy2018PMS}. 그 결과 명목상 동일한 부품이라도
기계적 성능이 서로 다를 수 있으므로, 오늘날의 인증은 witness 시편의 파괴
시험과 X-ray computed tomography
(XCT)~\cite{Seifi2016Qualification,Seifi2017Standardization,duPlessis2018XCT}에
크게 의존하며, 이러한 부품 단위 검사 부담은 산업적 확산을 가로막는 주요 비용
장벽으로 널리 지목되고 있다.

레이어 단위 in-situ 이미징은 이러한 부담에 대응할 수 있는 가장 폭넓게 이용
가능한 수단이다. 산업용 L-PBF 장비는 레이어마다 두 번---레이저 노광 직후 한
번과 분말 도포 직후 한 번---분말 베드를 촬영할 수 있으므로, 모든 빌드는
사실상 추가 비용 없이 자기 자신의 제조 과정에 대한 공간적으로 정합된 시각
기록을 남긴다~\cite{Scime2018Anomaly,Gobert2018InProcess,DebRoy2021NatRevMater,Mahmoud2021LPBFreview}.
이 기록을 활용하는 여러 방법 가운데 분말 베드 이상 영역의 픽셀 단위 의미
분할(semantic segmentation)은 두 가지 이유에서 특히 가치가 있다. 첫째, 공정
제어 측면에서, 픽셀 단위 결함 맵은 도포 및 용융 이상을 빌드 평면과 적층
방향을 따라 국소화하여 레이어 수준의 개입을 지원하며, 그 클래스별 구조는
리코터 줄무늬(streaking)와 같은 장비 상태 신호를 과용융·불완전 용융과 같은
에너지 밀도 신호로부터 분리해 낸다~\cite{Scime2020DSCNN}. 둘째, 물성 예측
측면에서, Scime 등은 세그멘테이션 맵의 슈퍼복셀(supervoxel) 집계량을 다른
in-situ 피처와 결합하면 6{,}299개 시편에 걸쳐 국소 인장 특성을 예측할 수
있음을 보였다~\cite{VPPM_Baseline,ORNL_CoRegistered_Dataset}. 그러나
세그멘테이션 정보원은 하류의 모든 것을 제한한다.
Section~\ref{sec:results}에서 보이듯이, 그 외 모든 조건이 동일한 물성 예측
파이프라인에 배포된 dynamic segmentation convolutional neural network
(DSCNN)~\cite{Scime2020DSCNN}가 데이터셋에 제공한 결함 채널만을 입력하면
인장강도(ultimate tensile strength, UTS) 예측의 평균 제곱근
오차(root-mean-square error, RMSE)가 단순(naive) 전역 평균 예측기보다도 나쁜
수준으로 붕괴하는 반면, 동일한 파이프라인에 본 논문이 제안하는 모델의 결함
맵을 입력하면 그 오차가 단순 예측기 대비 절반 가까이 줄어든다.

더 나은 분말 베드 세그멘테이션을 가로막는 중심 장애물은 어노테이션 부족이다.
수 메가픽셀 규모의 레이어 이미지를 전문가가 픽셀 단위로 라벨링하는 작업은
느리고, 그 결과로 공개된 코퍼스의 규모는 작다: Peregrine v2021-03 릴리스는
20개의 전문가 어노테이션 L-PBF 레이어를 포함하며~\cite{Peregrine_v2021_03},
v2022-10 릴리스의 여섯 개 장비--소재 조합은 각각 5개에서 26개의 어노테이션
레이어를 제공한다~\cite{Peregrine_v2022_10}. 본 연구에서 사용한 여덟 개
소스를 모두 합해도 사람이 어노테이션한 레이어는 125개에 불과하다. 한편
이미지의 외관은 장비, 광학계, 조명, 소재에 따라 크게 달라지므로, 한 플랫폼의
수십 개 레이어만으로 처음부터 학습된 세그멘테이션 네트워크는 심각한 과적합
위험과 함께, 다른 플랫폼으로 전이될 때 고전적인 데이터셋 시프트(dataset
shift) 페널티를 동시에 떠안는다~\cite{Quinonero2009DatasetShift,Torralba2011DatasetBias,Pan2010TransferLearning}.
이에 대한 사실상의 표준적 대응---DSCNN--Peregrine
계열~\cite{Scime2020DSCNN}과 같이 장비마다 별도의 지도학습 모델을 학습하는
방식---은 라벨링 비용을 장비 대수만큼 배가시킨다.

본 논문은 고정된(frozen) 자기지도 비전 파운데이션 모델(foundation model) 위에
구축된 분말 베드 결함 세그멘테이션 모델인 \textbf{DefectSegAM}을
제안한다(Fig.~\ref{fig:architecture}). 레지스터 토큰(register token)을 갖는
DINOv2 ViT-B/14 인코더~\cite{Oquab2024DINOv2,Darcet2024Registers}는 전체가
고정된 채 유지되며, 네 개의 트랜스포머 깊이에서 추출된 피처는 레이어당 두
장의 이미지(용융 직후와 도포 직후)에 걸쳐 융합되고---이때 베드의 용융 유발
변화를 인코딩하는 명시적인 이미지 간 차분(difference) 항을 포함한다---DPT
방식의 다중 스케일 디코더~\cite{Ranftl2021DPT}로 디코딩된다. 고정된 86~M개의
백본 파라미터 위에서 디코더와 헤드의 14.6~M개 파라미터만이 학습된다. 학습은
라벨 효율적인 두 단계로 진행된다: Stage~1에서는 ORNL co-registered
데이터셋~\cite{ORNL_CoRegistered_Dataset}의 다섯 개 빌드 중 네 개에서 추출한
12{,}607개의 미어노테이션 양산 레이어에 대해, 배포된 DSCNN의 자체 예측을
비용이 들지 않는 의사 라벨(pseudo-label)로 사용하는 지식 증류(knowledge
distillation)로 디코더를 사전학습하고, Stage~2에서는 여섯 개의 L-PBF 소스와
두 개의 바인더 젯팅(binder jetting) 소스에서 취합되어 통일된 11-클래스
온톨로지(Table~\ref{tab:classes})로 재매핑된 사람 어노테이션 레이어에 대해
미세조정한다~\cite{Peregrine_v2021_03,Peregrine_v2022_10}. ORNL 코퍼스의
극단적인 클래스 불균형(Section~\ref{sec:method-problem})은
역빈도(inverse-frequency) 클래스 인지 샘플링, Copy-Paste
증강~\cite{Ghiasi2021CopyPaste}, focal
loss~\cite{Lin2017FocalLoss}로 대응한다. 아울러 향후 어노테이션을 위한
외형 기준(appearance-based) 결함 분류 체계를
제안한다(Section~\ref{sec:method-taxonomy}).

평가 설계에서는, ORNL 양산 빌드의 밀집 라벨이 DSCNN의 예측이어서 해당 빌드의
지표가 교사(teacher) 일치도에 그치므로(Section~\ref{sec:setup-metrics}), 정확도
주장은 모든 학습에서 제외된(held-out) 사람 ground truth만을 기준으로 한다:
(i) 외부 2025 Peregrine held-out
세트~\cite{Peregrine_v2025_TODO}(Section~\ref{sec:setup-heldout-peregrine}),
그리고 (ii) 장비 간(cross-machine) 일반화를 검증하기 위해 저자들이 현재
수집 및 픽셀 라벨링을 진행하고 있는 Veltz M160 및 M280 장비와 EOS PBF
장비의 리코팅 이미지(Section~\ref{sec:setup-heldout-crossmachine})가
그것이다. DefectSegAM과 동일한 데이터로 학습한
DSCNN~\cite{Scime2020DSCNN}의 충실한 재구현을 주요 베이스라인으로 사용한다. \TODO{held-out 사람 GT 실험과 장비 간
실험(Sections~\ref{sec:results-heldout}
및~\ref{sec:results-crossmachine})은 계획되어
있으나 아직 수행되지 않았다; 결과가 나오는 대로 이 문단에 대표 수치를 삽입할
것.}

본 연구의 기여는 다음과 같다.
\begin{enumerate}
  \item \textbf{아키텍처.} 각 레이어의 도포 직후 이미지와 용융 직후 이미지에
        공유되는, 레지스터 토큰을 갖는 고정 DINOv2 ViT-B/14 인코더를,
        명시적 차분 항을 포함하는 단계별(per-stage) 이중 이미지 융합 및 DPT
        방식 다중 스케일 디코더와 결합한 DefectSegAM을 제안한다. 고정 백본
        위에서 디코더와 헤드만이 학습되므로, 시각 표현은 플랫폼마다 다시
        학습되는 대신 장비와 소재 전반에 걸쳐 상각(amortize)된다.
  \item \textbf{라벨 효율적 2단계 학습.} 배포된 양산 모델의 의사 라벨
        레이어로부터의 지식 증류와, 통일된 11-클래스 온톨로지 아래 여덟 개의
        이종 장비--소재 소스에서 취합한 소규모 사람 어노테이션 레이어에 대한
        미세조정을 결합하고, 여러 자릿수에 걸쳐
        차이 나는 클래스별 픽셀 수에는 클래스 인지 샘플링, Copy-Paste 증강,
        focal loss로 대응한다.
  \item \textbf{Held-out 사람 GT 평가 프로토콜.} 사람 GT만을 정확도의
        기준으로 삼는 평가 스위트를 명시한다: 학습에 전혀 사용되지 않은
        외부 2025 Peregrine 사람 어노테이션 벤치마크, Veltz M160/M280 및
        EOS PBF 장비의 장비 간 테스트 이미지(수집 및 라벨링 진행 중),
        그리고 동일한 데이터로 충실히 재구현·학습된 DSCNN 베이스라인이
        그것이다.
        \TODO{프로토콜은 정의되었으나 실험은
        미수행(Section~\ref{sec:setup}).}
  \item \textbf{하류 유용성 사례 연구.} 이미 발표된 인장 특성 예측
        파이프라인 전체를 고정한 채 세그멘테이션 정보원만을 교체하는 실험을
        통해, 데이터셋이 제공하는 DSCNN 맵은 단순 전역 평균 예측기 아래로
        붕괴하는 반면(UTS RMSE 70.7 대 68.4~MPa), DefectSegAM 맵은 단순
        베이스라인 대비 45.2\% 낮은 37.5~MPa에 도달하며 네 가지 인장 특성
        모두에서 개선을 보임을 확인한다(Section~\ref{sec:results-downstream}).
  \item \textbf{외형 기준 결함 분류 체계(부수적 제안).} 실험과 별도로,
        Peregrine 어휘에 혼재된 원인 기준과 외형 기준을 분리하여 레이어
        이미지 특징을 판정 기준으로 정의한 결함 8종(Under Melting 처리 확정 전
        기준)과 배경 2종의 체계를 향후 어노테이션을 위해 제안하고, 기존
        라벨 공간 및 ASTM E3353-22 용어와의
        대응(Table~\ref{tab:taxonomy})과 데이터 가용성을
        정리한다(Section~\ref{sec:method-taxonomy}).
        \TODO{제안 체계로 재어노테이션한 뒤 재학습·재평가할
        것(Section~\ref{sec:method-taxonomy}).}
\end{enumerate}

본 논문의 나머지 부분은 다음과 같이 구성된다. Section~\ref{sec:related}에서는
분말 베드 모니터링, ex-situ 라벨 기반 결함 검출, AM 분야의 파운데이션 모델,
결함-물성 모델링에 관한 관련 연구를 검토하고 연구 공백을 종합한다.
Section~\ref{sec:method}에서는 문제 설정과 실험용 11-클래스 온톨로지,
제안하는 외형 기준 결함 분류 체계, DefectSegAM의 아키텍처, 2단계 학습
레시피, 클래스 불균형 대응 전략, 그리고 증강·추론·후처리 파이프라인을
차례로 상술한다. Section~\ref{sec:setup}에서는 데이터셋, held-out 사람 GT
평가 프로토콜, 베이스라인, 평가 지표, 구현 세부사항과 다운스트림 물성
예측 설정을 기술한다. Section~\ref{sec:results}에서는 세그멘테이션 및 하류
물성 예측 결과를 보고하고, Section~\ref{sec:discussion}에서는 한계와 시사점을
논의하며, Section~\ref{sec:conclusion}에서 결론을 맺는다.
